\documentclass{exam}
\usepackage{math-macros}

\department{}
\school{}
\teacher{Vũ Vương An}
\topic{Ứng dụng đạo hàm để khảo sát và vẽ đồ thị hàm số}
\examtitle{PHIẾU BÀI TẬP}
\subject{Toán}
\grade{12}
\duration{20 phút}
\examdate{17/09/2026}
\setquestiontotal{ 10 }

\answerkeyfalse


\begin{document}
\makeheader


\begin{question}
Hàm số $y = x^4 - 2x^2$ có bao nhiêu điểm cực tiểu?



\choicegrid{ $0$ }{ $1$ }{ $2$ }{ $3$ }


\end{question}




\begin{question}
Cho hàm số $y = x^3 - 3x^2 + 2$. Hoành độ điểm cực đại của đồ thị hàm số là



\choicegrid{ $-1$ }{ $0$ }{ $1$ }{ $2$ }


\end{question}




\begin{question}
Hàm số $y = -x^3 + 3x$ có bao nhiêu điểm cực trị?



\choicegrid{ $0$ }{ $1$ }{ $2$ }{ $3$ }


\end{question}




\begin{question}
Hàm số $y = x^3 + 3x^2$ nghịch biến trên khoảng nào dưới đây?



\choicegrid{ $(-\infty; -2)$ }{ $(-2; 0)$ }{ $(0; +\infty)$ }{ $(-1; 1)$ }


\end{question}




\begin{question}
Giá trị nhỏ nhất của hàm số $y = x^3 - 3x + 2$ trên đoạn $[0; 2]$ bằng



\choicegrid{ $-2$ }{ $-1$ }{ $0$ }{ $1$ }


\end{question}




\begin{question}
Giá trị lớn nhất của hàm số $y = x^3 - 3x + 1$ trên đoạn $[-1; 2]$ bằng



\choicegrid{ $1$ }{ $2$ }{ $3$ }{ $5$ }


\end{question}




\begin{question}
Hàm số $y = x^3 - 3x$ đồng biến trên khoảng nào dưới đây?



\choicegrid{ $(-1; 1)$ }{ $(0; 2)$ }{ $(1; +\infty)$ }{ $(-2; 0)$ }


\end{question}




\begin{question}
Một chất điểm chuyển động trên trục $Ox$ với vận tốc $v(t) = 3t^2 - 12t + 9$ (m/s), $t$ tính bằng giây. Thời điểm chất điểm đứng yên lần đầu tiên là



\choicegrid{ $1$ s }{ $2$ s }{ $3$ s }{ $4$ s }


\end{question}




\begin{question}
Tìm $m$ để hàm số $y = x^3 - 3x^2 + m - 1$ đạt cực đại bằng $0$.



\choicegrid{ $0$ }{ $1$ }{ $2$ }{ $3$ }


\end{question}




\begin{question}
Đồ thị hàm số $y = x^3 - 3x$ cắt đường thẳng $y = m$ tại ba điểm phân biệt khi



\choicegrid{ $m = 2$ }{ $m = -2$ }{ $m = 0$ }{ $m = 3$ }


\end{question}





\end{document}
